\subsection{Orienting the fresh two-object search by a nonzero restriction}

\begin{lemma}[Proved: a fixed-point-free positive source orientation]
\label{CP2:A-two-object-positive-orientation}
Suppose a two-object collage $C$ has classifying space
of type $\CP^2$. After possibly taking the opposite
category, its source corner $M$ has nonzero second
rational homology and its connecting right $M$ act has
no global fixed state. In fact a generator of
$H^2(BC;\mathbb Q)$ restricts nontrivially to that source
corner.
\end{lemma}
\begin{proof}
Let $\alpha$ be the rational degree-two generator.
Its square is nonzero. By the relative-cup lemma,
$\alpha$ cannot restrict to zero on both corners.
Choose a corner on which its restriction is nonzero,
and if needed take opposition to make it the source.
Opposition preserves the homotopy type and the
nonvanishing of that restriction. A global fixed state
of the resulting right act would supply a natural
transformation from the source inclusion to the constant
functor at the other object, making its restriction zero.
This is impossible.
\end{proof}

\begin{corollary}[Proved: the joint fixed-point-free eighteen-arrow scope]
\label{CP2:A-two-object-eighteen-oriented-scope}
To close the fresh two-object eighteen-arrow lower-bound
scope, it suffices to check the twenty-five ordered
positive-source profiles of
Proposition~\ref{CP2:A-two-object-eighteen-scope}, using
only right acts without a global fixed state, provided
all otherwise required opposite-corner tables are included.
In particular, $P_5$ targets in profiles with source order
six through nine must be included under this orientation
convention. They cannot be omitted merely by citing the
oppositely ordered profile with a fixed-point-free $P_5$
right act.
\end{corollary}
\begin{proof}
Choose the orientation from the lemma before applying
the complete eighteen-arrow cardinality and analytic
reductions. The ordinary positive-source catalogs include
both ordinary opposition classes, so the oriented source
is present. The remaining ordered size triple is among
the same twenty-five profiles. Every retained analytic
input reduction continues to apply to this subfamily.
Conversely, the finite checks must now include every
allowed opposite table: exchanging corners again might
lose the specified nonzero restriction and fixed-point-free
right action. This explains the explicit $P_5$ target
requirement. Once all these ordered scopes are joined,
the lemma supplies a checked orientation for every
possible projective-plane collage.

For minimal-band right-act normal forms, a state fixed
by the whole band is automatically fixed by the entire
source monoid: $sm=(se)m=s(em)=s$ for a band idempotent
$e$, because $em$ remains in the band. Thus the literal
fixed-state metadata of the complete band catalogs
tests this condition before assigning outside actors.
For a nontrivial minimal group the same argument uses
a state fixed by every element of the minimal ideal.
\end{proof}

\begin{remark*}[Scope accounting]
The new convention is a joint orientation reduction.
An isolated fixed-point-free profile receipt does not
exclude all categories of that ordered profile by itself.
It becomes an unrestricted two-object exclusion only
after the complete ordered profile join and its opposite
dependencies are checked. Earlier complete seventeen-arrow
receipts and the unpruned mixed-five eighteen-arrow run
are unchanged.
\end{remark*}
